\documentclass[11pt,letterpaper]{article}
\usepackage[T1]{fontenc}
\usepackage{lmodern,amsmath,amssymb,amsthm,mathtools}
\usepackage[margin=0.82in]{geometry}
\usepackage{microtype,booktabs,longtable,array,enumitem,xcolor,fancyhdr,xurl}
\usepackage[colorlinks=true,linkcolor=blue!45!black,urlcolor=blue!45!black,citecolor=blue!45!black]{hyperref}
\usepackage{listings}
\hypersetup{pdftitle={Sparse syntax and weighted Pell transport for Diophantine consistency},pdfauthor={Research prepared with ChatGPT for Vladimir Reshetnikov and ProveIt}}
\newtheorem{theorem}{Theorem}[section]
\newtheorem{lemma}[theorem]{Lemma}
\newtheorem{proposition}[theorem]{Proposition}
\newtheorem{corollary}[theorem]{Corollary}
\theoremstyle{remark}\newtheorem{remark}[theorem]{Remark}
\newcommand{\Z}{\mathbb Z}\newcommand{\N}{\mathbb N}
\newcommand{\Con}{\operatorname{Con}}\newcommand{\RefZ}{\operatorname{Ref}}
\newcommand{\code}[1]{\nolinkurl{#1}}
\setlength{\parskip}{4pt}\setlength{\parindent}{0pt}
\setlength{\emergencystretch}{2em}
\setlist{itemsep=2pt,topsep=4pt}
\allowdisplaybreaks[2]
\pagestyle{fancy}\fancyhf{}
\fancyhead[L]{\small ProveIt research memorandum}
\fancyhead[R]{\small Sparse syntax and Pell transport}
\fancyfoot[C]{\thepage}
\setlength{\headheight}{14pt}
\lstset{basicstyle=\ttfamily\small,breaklines=true,columns=fullflexible,keepspaces=true}
% Exact endpoint metrics; checked against direct208/metrics.json by verify_all.py.
\newcommand{\DirectVariables}{208}
\newcommand{\DirectRows}{148}
\newcommand{\DirectLastRow}{147}
\newcommand{\DirectLastVariable}{207}
\newcommand{\DirectHeightMath}{\ensuremath{2^{44}}}
\newcommand{\DirectMonomials}{5,987}


\title{Sparse syntax and weighted Pell transport\\
\large Explicit Diophantine equations for the consistency of ZFC}
\author{Prepared for Vladimir Reshetnikov and the ProveIt project\\
\small Mathematical development and independent internal reviews assisted by ChatGPT}
\date{26 October 2026}

\begin{document}
\maketitle
\begin{abstract}
We give two explicit improvements to the inspected ProveIt constructions.
The direct construction uses a sparse, non-surjective coding of its numerical
proof calculus. A finite-marker substitution theorem proves that the unused
numeric codes do not change finite refutability. Combining this coding with
a signed affine horizon and divisibility-forced Pell growth produces a
quartic in \DirectVariables{} unrestricted integer variables, specified by
\DirectRows{} quadratic residuals. The separate compact construction gives a
polynomial of exact degree 190 in ten strictly positive integer variables,
with 120 binary arithmetic operations and a completely specified fixed
coefficient tuple. Its new weight $b^2UE$ permits a whole-product Pell center
in a smaller arithmetic chart, including an auxiliary index previously
excluded in the other chart. For each construction, ordinary proofs of the
new steps compose with explicitly identified inherited calibrations to prove
in ZFC that absence of solutions is equivalent to syntactic $\Con(\mathrm{ZFC})$.
All coefficients, evaluators, exact finite checks, source bindings, and
editable article sources accompany the paper. Neither global minimality nor
a new end-to-end proof-assistant verification is claimed.
\end{abstract}

\tableofcontents
\newpage

\section{Results, comparison, and the precise logical claim}

Fix the effective ZFC presentation used in the repository's retained
calibration: its axiom predicate is \code{ZFCax_s}, its finite derivation
relation is \code{BProv}, and its false formula is \code{Form.fBot}.
Write $\RefZ_{\mathrm{ZFC}}$ for the arithmetic assertion that a finite
refutation in this presentation exists, and put
$\Con(\mathrm{ZFC})=\neg\RefZ_{\mathrm{ZFC}}$.

Our comparison snapshot is the public ProveIt commit
\code{d123610bc98fef2006bff1c9466bc3f1221003e7}. We inspected the supplied
Hilbert Tenth Problem directory, the numbered consistency constructions
through package 178, and the seven archives then in \code{docs/incoming}.
The incoming Pell Growth 210 archive supplies the immediate direct parent\cite{growth210}.
Packages 172 and 178, together with the incoming affine-horizon archive,
supply the compact comparison and its two principal ingredients\cite{free172,center178}.
These are comparisons against identified files; they are not a proof of
priority or an exhaustive worldwide search.

\begin{center}
\small
\begin{tabular}{@{}lrrrr@{}}
\toprule
Construction and domain&Variables&Degree&Rows or operations&Height\\
\midrule
Incoming Pell Growth, $\Z$&210&4&150 rows&$2^{40}$\\
Affine-horizon companion here, $\Z$&209&4&149 rows&$2^{44}$\\
\textbf{Sparse syntax here, $\Z$}&\textbf{\DirectVariables}&\textbf{4}&\textbf{\DirectRows{} rows}&\DirectHeightMath\\
\addlinespace
Incoming free chart, $\Z_{>0}$&10&206&90 operations&fixed tuple\\
Package 178 whole center, $\Z_{>0}$&10&208&116 operations&fixed tuple\\
\textbf{Weighted transport here, $\Z_{>0}$}&\textbf{10}&\textbf{190}&\textbf{120 operations}&fixed tuple\\
\bottomrule
\end{tabular}
\end{center}
Height for the direct constructions means the largest absolute coefficient after the
single quartic has been expanded and like monomials collected. The new
direct quartic has \DirectMonomials{} collected monomials; its parent has
5,282. The companion has 5,486. The displayed tradeoffs are intentional:
dimension, degree, expanded height, support, and arithmetic-circuit length
are different measures of simplicity.

\begin{theorem}[Delivered equations]\label{thm:main}
Let $R_0,\ldots,R_{\DirectLastRow}$ be the literal quadratic residuals in
Appendix~\ref{app:residuals}, and define
\[
 Q_{\DirectVariables}(x_0,\ldots,x_{\DirectLastVariable})
       =\sum_{j=0}^{\DirectLastRow}R_j(x_0,\ldots,x_{\DirectLastVariable})^2.
\]
Let $P$ be the ten-variable polynomial in
Section~\ref{sec:weighted-positive}, using the four fixed integers completely
defined in Appendix~\ref{app:coefficients}. Then ordinary mathematics
formalizable in ZFC proves
\begin{align*}
 \exists\boldsymbol x\in\Z^{\DirectVariables}\;
       Q_{\DirectVariables}(\boldsymbol x)=0
       &\quad\longleftrightarrow\quad\RefZ_{\mathrm{ZFC}},\\
 \exists\boldsymbol u\in\Z_{>0}^{10}\;P(\boldsymbol u)=0
       &\quad\longleftrightarrow\quad\RefZ_{\mathrm{ZFC}}.
\end{align*}
Their negations are therefore ZFC-provably equivalent to
$\Con(\mathrm{ZFC})$. The first equation has exact degree four and the
second has exact degree 190.
\end{theorem}

The quantifiers in the first assertion range independently over all
integers. The second assertion has a strictly positive domain. Its named
registers are abbreviations for polynomials in ten inputs, so they are not
extra existential witnesses. The fixed coefficient recipe uses finite
sums, digit extraction, and fixed powers only to name four particular
integers. No variable exponentiation or floor occurs in either delivered
polynomial.

The direct route evaluates the numerical calculus C222 itself and uses its
identified C222/ZF/ZFC calibration. It does not simulate K215 or specialize
Jones's universal polynomial. The compact route retains its existing
native/K215 ancestry. Allowing alternative sources is useful, but replacing
an already effective source is not automatically a simplification.

\paragraph{Proof status.}
The new results are ordinary mathematical arguments, supported by exact
coefficient checks. We have not emitted an object-language ZFC derivation
or rebuilt a complete transitive Lean/Rocq dependency closure. The retained
logical endpoints are \code{KEndpoint.c222_iff_zf_refutation} and
\code{ConstructibleChoice.zf_refutation_iff_zfc_refutation} on the direct
side, and the explicitly named whole-tuple native/refutation calibration on
the compact side. Finite checks authenticate formulas and transformations;
they do not prove unbounded Pell or logical theorems by sampling. The result
assumes neither consistency nor inconsistency and proves neither one.

\section{Sparse syntax: a finite proof does not need a surjective coding}
\label{sec:sparse}

The earlier square-shell coding makes every natural number a constructor
value. That property simplifies global recoding, but it is stronger than
what is needed to preserve the existence of a finite refutation. Define
\begin{equation}\label{eq:sparse-constructors}
 \begin{aligned}
 \pi(x,y)&=(x+y)^2+x,\\
 I(x,y)&=4\pi(x,y),& M(x,y)&=4\pi(x,y)+1,\\
 A(x,y)&=4\pi(x,y)+2,& E(x,y)&=4\pi(x,y)+3.
 \end{aligned}
\end{equation}
Here $I,M,A,E$ name the implication, membership, quantifier, and equality
constructors of the same numerical rule schemata. They are polynomials,
not syntactic predicates. Their arguments may be any natural numbers.

For $x+y=s$, $\pi(x,y)$ runs injectively through
$[s^2,s^2+s]$. Different shells have disjoint intervals. The four
constructor ranges are disjoint by their residues modulo four. The first
unused values are $12,13,14,15$, because the quotient 3 is not in the range
of $\pi$. Every nonzero constructor value is greater than both arguments.
The only zero case is $I(0,0)=0$, and $M(0,0)=1$.

\subsection{Normalized trees and finite markers}

Decode a numeral recursively: zero is a distinguished leaf; a nonzero
constructor value is its uniquely determined constructor node; a number
outside all constructor ranges is a new atomic leaf. Strict descent of
the arguments makes this recursion terminate. Trees are normalized by
identifying $I(0,0)$ with zero. Apart from this convention, injectivity
and disjointness give unique trees.

The reference Cantor constructors are
\[
 I_o(x,y)=2(x+y)(x+y+1)+4y,\qquad
 M_o=I_o+1,\quad A_o=I_o+2,\quad E_o=I_o+3.
\]
They partition $\N$ and have the same descent and zero convention. Every
reference numeral therefore has a normalized tree built from zero alone.

\begin{lemma}[Finite-marker substitution]\label{lem:markers}
Let $\mathcal S$ be a finite, subtree-closed set of normalized sparse-code
trees. There is a substitution of reference-code numerals for the extra
atomic leaves that is injective on $\mathcal S$, fixes 0 and 1, and
preserves every constructor occurrence represented in $\mathcal S$.
\end{lemma}
\begin{proof}
Let $D$ be the maximum height in $\mathcal S$ and let there be $k$ distinct
extra leaves. If $k=0$, evaluate the same trees in the reference coding.
Otherwise put $L=\max(D,k)$ and define
\[
 W_0=0,\qquad W_{r+1}=A_o(0,W_r),\qquad
 m_j=I_o(W_L,W_j)\quad(0\le j<k).
\]
The markers $m_j$ are distinct and all have height $L+1>D$. Substitute
the $j$th marker for the $j$th extra leaf.

A tree without extra leaves has height at most $D$. The image of one
extra leaf has height $L+1$. The image of a constructor tree containing
an extra leaf has height greater than $L+1$, with height determined by
the deepest occurrence of an extra leaf. Thus a marker cannot collide
with a constant-only tree or a nontrivial tree containing markers.
If two remaining nonzero constructor trees have equal images, disjoint
constructor ranges identify the root symbol and injectivity identifies
both child images. Induction then identifies the original trees. The
normalization $I(0,0)=0$ is imposed on both sides, so it introduces no
new collision. Substitution fixes zero and $1=M(0,0)$ and commutes with
the constructors by definition.
\end{proof}

\begin{proposition}[Finite refutation equivalence]\label{prop:recoding}
The sparse C222 calculus derives 1 if and only if the reference C222
calculus derives 1. ZFC proves the equivalence by effective finite proof
transformations.
\end{proposition}
\begin{proof}
For a finite sparse derivation, collect every parameter, premise,
conclusion, and intermediate value of every instantiated constructor
term. Decode these values and include all subtrees. Lemma~\ref{lem:markers}
transfers each instantiated rule to a reference rule. It preserves truth
at zero, modus ponens, generalization, and all constructor-term axioms.
Its injectivity preserves the guarded bundle's two disequalities
$c\ne a$ and $c\ne b$. The concluding numeral remains 1.

Conversely, evaluate the unique zero-generated tree of each reference
numeral with the sparse constructors. Uniqueness of sparse decoding
makes this a global injection on reference numerals. It commutes with
all constructors, fixes 0 and 1, and transports every finite derivation
and its guards to the sparse calculus.

The arithmetic bundle guard uses the existential identity
$((c-a)(c-b))^2=(d+1)^2$. After translating a guarded instance, choose
$d=|(c-a)(c-b)|-1\ge0$. One must not assume that a numeral recoding
commutes with addition or preserves the old $d$. In other rule branches,
$d$ is simply a schema parameter and is translated in the ordinary way.
These are alternative local branches, so the choices do not conflict.
\end{proof}

There is generally no global injective homomorphism in the reverse
direction fixing zero: the reference constructor closure of zero already
exhausts the reference numerals. The dependence on the finite derivation
is essential. It is also sufficient, because inconsistency is the
existence of one finite refutation. This is the principal conceptual
change in the direct construction.

\section{Signed horizons and a compatible Pell growth certificate}
\label{sec:growth}

For an integer $a\ge2$, use the canonical Pell notation
\[
 (a+\sqrt{a^2-1})^n=X_n(a)+Y_n(a)\sqrt{a^2-1}
 \qquad(n\in\Z).
\]
In the compact construction's inherited notation, $\psi_a(n)=Y_n(a)$.

The arithmetic source uses a horizon $t$, a growth value $F$, a beta
spacing $y$, a beta code $x$, and an even CRT index representative $C$.
They are ordinary integer witnesses before their signs have been proved.
Write $d=2^{19}$. Replace the old positive square slice by the affine
slice
\begin{equation}\label{eq:affine-horizon}
 t=2+4z,\qquad z\in\Z.
\end{equation}
For the sparse syntax use the three quadratic growth equations
\begin{equation}\label{eq:sparse-growth}
 Z_g=tF,\qquad X_g^2-64Z_g^2-16Z_gF=1,
 \qquad F=dt(2u+1).
\end{equation}
The changed center is $8t+1$. Keeping the old center $4t+1$ without
updating the label bound would leave a proof gap; the literal final
equation contains the new coefficients 64 and 16.

\subsection{Recover the sign before classifying the growth Pell equation}

The retained source has a common square frame and a largest-index
certificate for $C$. Its inherited theorem recovers each principal Pell
pair at its signed nominal index before the front variables are oriented.
For a requested base with absolute value at least two, its signed power
lemma then forces the exponent to be nonnegative and the output to be
the exact power. The actual frame is retained with every requested base,
output, and nominal exponent included.

Equation~\eqref{eq:affine-horizon} gives even $t\ne0$, $|t|\ge2$.
The last row of~\eqref{eq:sparse-growth}, without using its norm row,
already gives
\[
 F\ne0,\qquad F\equiv0\pmod2,\qquad |F|\ge d|t|\ge2^{20}.
\]
In particular the existing square frame exceeds its required threshold
258, because it includes $F^2$.

The affine radix is $D_r=2(C+F+1)$. Since $C$ and $F$ are even,
$|D_r|\ge2$. The retained request $q=D_r^x$ gives $x\ge0$.
The beta equations are
\begin{equation}\label{eq:beta-endpoints}
 y=F(2v+1),\qquad
 x=F+1+(y+1)r_0,\qquad
 x=1+[1+(t+1)y]r_t.
\end{equation}
They imply $|y|\ge|F|$ and
\[
 |1+(t+1)y|\ge|t+1|\,|y|-1\ge|F|-1>1.
\]
Consequently $x=0$ would express $-1$ as a multiple of a nonunit,
which is impossible. Thus $x\ge1$, $|q|\ge2$, and the existing request
at base $q$ with exponent $t$ forces $t\ge0$. The affine slice excludes
zero, so $t\ge2$.

This bootstrap uses the elementary divisibility lower bound before the
standalone Pell center has been oriented. It never assumes a growth
estimate at a negative center. It is the new composition of the signed
horizon and the incoming divisibility-growth construction.

\subsection{Label bound and growth}

For $0\le r,s\le M$, each sparse constructor satisfies
\[
 C(r,s)+1\le16M^2+4M+4\le16(M+1)^2.
\]
Every leaf of a local constructor term is below $t$ or is 0 or 1.
Induction on depth gives the bound
$16^{2^h-1}t^{2^h}$ for a term's value plus one. The deepest source
term has depth 19; the eight ordinary axiom terms have depth at most
five. Define
\[
 L=16^{d-1}t^d.
\]
Every required local formula label is less than $L$.

After the sign bootstrap, $a=8t+1\ge17$ and the norm row is
$X_g^2-(a^2-1)F^2=1$. Write $|F|=Y_n(a)$ with $n>0$ for its canonical
Pell ordinate. The retained elementary identities give
\[
 Y_n(a)\equiv n\pmod{a-1},\qquad
 v_2(Y_n(a))=v_2(n)\quad(a\text{ odd}).
\]
Since $8t\mid F$, one has $8t\mid n$. The exact valuation of
$F=dt(2u+1)$ then shows that $n/(8t)$ is divisible by $2^{16}$,
and hence $n\ge dt$. The recurrence growth estimate yields
\[
 |F|\ge(16t+1)^{dt-1}\ge(16t)^{dt-1}\ge L^t,
\]
because the ratio in the last comparison is $16^{t-1}/t\ge1$.
In particular $|F|>t$.

The retained requests at exponents $y+1$ and $F+t$, together with
\eqref{eq:beta-endpoints}, now recover
\[
 F>0,\quad y\ge F,\quad x\ge2+(t+1)y,\quad C>F+t,
 \quad D_r>2C,\quad q=D_r^x>0.
\]
The common bound $F\ge L^t$ is exactly what is needed by the inherited
finite-product separation and Horn-flow argument. It also gives
$F\ge t^{2t}\ge\operatorname{lcm}(1,\ldots,t)^2$, as required by
the generalized factorial-free CRT selector.

\subsection{An explicit new bound for the changed local polynomial}

The syntax change must also respect the modulus used to recover the
local row. Put $R_0=2+t+x+y+F$. Every decoded local atom is less than
$R_0$, and $R_0\ge16$. If two inputs are below $R_0^e$, a sparse
constructor is below $16R_0^{2e}\le R_0^{2e+1}$. Thus a depth-19
label is below $R_0^{2^{20}-1}$ and a depth-five label is below
$R_0^{63}$.

The eleven unsquared factors in the local axiom/inference product have
the following valid exponent bounds: 1 for the truth factor, 66 for
each of eight ordinary axiom factors, 10 for modus ponens, and 7 for
generalization. Their squared product is below $R_0^{1092}$. The
squared bundle factor is below $R_0^{2097156}$ and its squared guard
defect is below $R_0^{10}$. The joint factor is therefore below
$R_0^{2097157}$. Each of the two beta squares is below $R_0^{12}$.
Including them gives
\[
 \mathrm{Local}<R_0^{2098250}<R_0^{2^{23}}.
\]
For $t\ge2$, one has $L\ge2^{5d-4}$ and hence
$F\ge L^t\ge2^{10d-8}>2^{32}$. The recovered inequalities imply
$R_0<2x$ and $x>2^{32}$. The function
$x/\log_2(2x)$ is increasing beyond this threshold, and at $2^{32}$
one has $2^{23}\cdot33<2^{32}$. Hence
\[
 \mathrm{Local}<(2x)^{2^{23}}<2^x.
\]
This supplies a fresh bound for the changed constructors rather than
silently reusing the old coefficients' bound. All remaining low-digit
estimates use $D_r>2C$ and the same authenticated index cap.

\subsection{Cofinal completeness}

Translate a finite ZFC refutation to a finite sparse-C222 refutation by
Proposition~\ref{prop:recoding} and the inherited calibration. Choose
a sufficiently large positive horizon with
\[
 t\equiv2\pmod{2^{23}}.
\]
This progression is cofinal and lies in \eqref{eq:affine-horizon}.
It has $v_2(t)=1$ and its binary digit in position 20 is zero. Set
\[
 n_0=dt,\quad y_0=Y_{n_0}(8t+1),\quad
 a_0=X_{n_0}(8t+1).
\]
The congruence gives $t\mid y_0$, and the exact valuation gives
$v_2(y_0)=20$. Together these facts, including the odd part of $t$,
show that $b_0=y_0/(dt)$ is an odd integer. Also $a_0\equiv1\pmod4$.

The incoming proof establishes that $s\mapsto Y_s(a_0)$ permutes
residues modulo every power of two when $a_0\equiv1\pmod4$.
Choose $N>23$ with $2^N>t$ and an arbitrarily large positive odd $s$
such that
\[
 (t/2)b_0Y_s(a_0)\equiv1\pmod{2^{N-20}}.
\]
Composition gives a positive ordinate
\[
 F=Y_{n_0s}(8t+1)=y_0Y_s(a_0)
       \equiv2^{20}\pmod{2^N},
\]
Complete its witnesses with
$X_g=X_{n_0s}(8t+1)$, $Z_g=tF$, and $u=(F/(dt)-1)/2$.
The latter is integral because the quotient is odd, so all three growth
rows hold. The binary supports of $F$ and
$t$ are disjoint, so $\binom{F+t}{t}$ is odd. This is the required
parity for an even CRT anchor, not an optional heuristic property.

Choose the odd multiplier in $y=F(2v+1)$ to contain every odd prime at
most $t$ and to be sufficiently large for the finite pending states.
The beta moduli, indexed by $1\le i<j\le t+1$, are pairwise coprime:
a common prime of $1+iy$ and
$1+jy$ would divide $i-j$, hence be at most $t$, and would also divide
$y$, a contradiction. The subsequent choice order is the beta code
$x$, the selector and CRT moduli, the ordinary CRT representatives,
an arbitrarily large even representative $C$, the radix $D_r$, its
power $q$, and finally the shared Pell frame. No later choice determines
a parameter on which an earlier modulus depends.

The retained arithmetic completeness theorem then supplies every
remaining witness. The local syntax block is regenerated with its new
constructors, its labels are covered by $L$, and its modulus bound was
proved above. This establishes the sparse source's existence equivalence
with the same finite refutation predicate.

\section{From a semantic source to the literal quartic}
\label{sec:compilation}

The source change and the polynomial reductions perform different jobs.
The former changes the numeric proof coding and horizon; its justification
is the unbounded mathematics of Sections~\ref{sec:sparse}--\ref{sec:growth}.
The latter are exact finite transformations of the newly defined residuals.

The local compiler represents every required constructor polynomial by
quadratic multiplication equations. Its common-expression rule is slightly
stronger than syntactic equality. If two quadratic expressions have
proportional homogeneous quadratic parts, with integer ratio $c$, then
\[
 q_2=cp_2,\qquad w=p_2+p_1
 \quad\Longrightarrow\quad q_2+q_1=cw+(q_1-cp_1).
\]
Here the subscripts 2 and 1 indicate the quadratic and affine parts.
The right side is affine in the existing wire. It therefore saves a
wire without division or a new integrality condition. The constructor
terms and all retained beta, CRT, and Pell rows are separately bound
by exact source certificates.

The graft's boundary is especially small: only the joint bundle/guard
square and the square of the eleven-factor axis product are passed
from the local block to the retained local-row equation. The independent
audit checks all nine axiom constructor expressions, the two inference
factors, both boundary values, and the retained-row partition. A stale
unused definition of $FS$ can be deleted or projected away; it cannot
be counted as a necessary witness.

The affine horizon eliminates its own $t$ coordinate. The common-Hessian
projections and subsequent integer coordinate/row changes have exact
restoration certificates. Coordinate shears and translations preserve
integer solutions bijectively. An invertible change of residual basis
preserves their simultaneous zero set and therefore preserves the zero
set of their sum of squares, although it need not preserve the sum-of-squares
polynomial itself. All these distinctions are recorded in the artifacts.

The final definition is the complete residual list, not merely the
existence of an algorithm that could someday print it. The file
\code{direct208/equations/principal208/equation.json} declares each input and each integer
coefficient. Its adjacent expansion lists every coefficient of the
single quartic. Appendix~\ref{app:residuals} prints the residuals.

The simpler 209-variable companion uses the original square-shell
syntax and changes only the horizon. Its exact parent-to-child
substitution is $t=2+4z$, deleting the unique horizon row and the
coordinate $t$. Three threshold choices, 2, 8, and 16, are retained in
\code{affine209/profiles}; the minimum-2 version has the smallest
coefficient height among those three. This companion isolates the
horizon improvement from the new sparse-syntax theorem.

\section{Weighted square transport: ten positive variables of degree 190}
\label{sec:weighted-positive}

The next equation improves the degree of the inspected ten-positive-variable
family. Its exact profile is
\[
 (\text{variables},\text{degree},\text{binary operations})=(10,190,120).
\]
The operation split is 75 multiplications, 24 additions, and 21 subtractions;
squares and multiplications by constants are charged. This is a tradeoff
against the incoming centered free chart's profile $(10,206,90)$, and against
source178's scaled whole-product center $(10,208,116)$. No global minimum or
worldwide priority is asserted. The new step is the polynomial transport
weight $b^2UE$, which permits whole-product centering in the free chart
without excluding auxiliary Pell index three.

\subsection{The complete polynomial and its fixed data}

The independent variables are
\[
 J,F,\alpha,i,T,w,s,n,\eta,\zeta\in\mathbb Z_{>0}.
\]
The four fixed integers $(\beta,K_0,M_C,M_F)$ are the complete, unchanged
BFS943/radix67 tuple of the frozen source172 parent. The archive supplies all
four numerical payloads and reconstructs them from its complete finite
436-character recipe. The largest integer has 592,884 bits. This is the same
whole tuple throughout the arithmetic and logical calibration; its scalar
bounds alone do not specify the logical instance.

All letters on the following display are polynomial abbreviations, not
additional quantified variables:
\begin{align*}
 h&=\beta J,&q&=h+1,&C&=q-F-\alpha,\\
 X&=qw,&Y&=s,&\bar Y&=qY,\\
 k&=\eta+\zeta,&c&=kY+\eta,&A&=\bar Y(X+1)+1,\\
 \Delta&=A^2-1,&H&=4A-5,&v&=X\bar Y^2,\\
 \ell&=vk,&E&=1+\ell(\ell+k),&U&=ic^2,\\
 R&=[h(q-F)+\alpha](q^2-1)+(M_C+qM_F)J,&&&\\
 B&=(K_0+w)C-F,&a&=\Delta U^2E,&F_0&=1+a,\\
 G&=EF_0,&f&=1+2a,&Q&=4\Delta aF_0,\\
 b&=cT-R,&V&=fb-c,&K&=X(3c-2X)+2,\\
 e_6&=6(k-R-1)-\ell(k^2-4),&D&=6q^7,&M&=hDH.
\end{align*}
Define the following division-free center:
\begin{align*}
 e_2&=(k\bar Y)^2,&
 S&=(k\bar Y)\bigl[\ell+k-(k\bar Y)(\bar Y+1)\bigr],\\
 \theta&=U^2b,&
 R_c&=\theta(3E+e_2)-(b-c),\\
 S_c&=4E^2+3S^2+e_2^2,&
 \Omega&=U\bigl[2SR_c+A(4EV-\theta S_c)\bigr].
\end{align*}
Finally put
\begin{align}
 L&=b^2UE,\label{eq:weighted-L}\\
 Z_0&=hD\bigl[H(n-LB)+K\bigr]
        -H\bigl[DB+h(e_6-6q^4s)\bigr],\notag\\
 Z&=Z_0-M\Omega,\notag\\
 P&=Q\bigl(Z^2-16M^2GV^2\bigr)+16M^2G-Z^2.\label{eq:positive190}
\end{align}
The single equation is $P=0$. Every operation displayed is ordinary integer
addition, subtraction, or multiplication with fixed exponents. The literal
120-operation graph is supplied in \texttt{positive190/equation.json}.
There is no quotient, radical, variable exponent, factorial, sequence
predicate, or unknown coefficient inside this polynomial.

\paragraph{Consistency theorem.}
At the identified inherited whole-tuple calibration, ZFC proves
\[
 \bigl(\neg\exists J,F,\alpha,i,T,w,s,n,\eta,\zeta>0\;P=0\bigr)
 \quad\longleftrightarrow\quad\operatorname{Con}(\mathrm{ZFC}).
\]
Here consistency is finite syntactic non-refutability for the fixed effective
ZFC presentation of the parent. The result establishes an equivalence; it
does not establish either consistency or unsolvability separately.

\subsection{The inherited interface and its precise point of use}

The pinned source172 ordinary proof supplies the following named ingredients:
prime-valuation square extraction; adjacent-coefficient rational Pell descent;
signed all-index folding and normalized strong-index recovery; the full Lucas
expansion; the $\lambda$-uniform ratio and complete-product projection theorem;
and fresh native reconstruction followed by the whole-tuple syntactic
refutation calibration. The selected source text is retained as
\texttt{inherited/FREE\_CHART\_PROOF.md}, with its original source pins.
The new proof does not claim a freshly compiled transitive Lean/Rocq closure.

For proof notation only, let
\[
 t=\frac Bh,\quad \kappa=\frac KH,\quad
 \sigma=\frac{Y}{q^3},\quad e^*=\frac{e_6}{6},\quad
 \rho=\frac{e^*}{q^7},\quad z=\frac ZM,\quad W=1+hL.
\]
These are rational numbers until the appropriate divisibilities are proved.
Direct algebra gives
\begin{equation}
 z=n+\kappa+\sigma-Wt-\rho-\Omega.
 \label{eq:weighted-carrier}
\end{equation}
The elementary positive-domain bounds are
\[
 q\ge32,\quad X\ge q,\quad\bar Y\ge q,\quad A>q^2,
 \quad c\ge3,\quad U\ge c^2,\quad E>1,
 \quad a>c,\quad f>2c,\quad Q>c^2,\quad M>0.
\]
At a zero, use the exact identities
\[
 Q=4\Delta^2U^2G,\qquad \gcd(E,F_0)=1,
 \qquad (Q-1)Z^2=16M^2G(QV^2-1).
\]
The prime valuations show that $G=g^2$, $E=\tau^2$, and $F_0=f_0^2$ for
positive integers $g,\tau,f_0$; the added factor16 is a square. The retained
adjacent-coefficient descent applies to the initially rational
$y=|z|/(4g)$ and proves it is integral. Thus
\begin{equation}
 z=\varepsilon4gy,\quad\varepsilon\in\{-1,1\},\qquad
 QV^2-(Q-1)y^2=1.
 \label{eq:weighted-root}
\end{equation}
These arguments depend only on the final norm equation, not the carrier.
Signed folding and the first index bound then give, before any assumption on
the signs of $B,C,R,b,V$,
\begin{gather}
 c=\psi_A(p),\quad p\ge3\text{ odd},\qquad
 c\ge4\Delta+3>4q^4,\quad K>0,\notag\\
 b\ne0,\quad\operatorname{sgn}V=\operatorname{sgn}b,
 \quad y<\tfrac32|V|,\quad |V|<\tfrac32 f|b|,
 \quad |b|>4\Delta f-1.\label{eq:weighted-inherited-bounds}
\end{gather}
These bounds explicitly allow auxiliary index three. Since the final
quantities are integers, $|b|\ge4\Delta f$. Also
\[
 G=E+\Delta(UE)^2\le A^2(UE)^2,
 \qquad g\le AUE.
\]

\subsection{The center bound in the changed chart}

The whole-product center of source178 was developed in a different chart.
Its polynomial identity is valid here with its letter $Y$ replaced by
$\bar Y=qY$. The bound must be checked before recovering the sign of $B$.
Temporarily let $u$ be a formal indeterminate, and write
\begin{align*}
 E(u)&=e_2u^2+e_1u+e_0,\\
 e_1&=k^2\bar Y(1-2\bar Y-2\bar Y^2),\\
 e_0&=1+k^2\bar Y(\bar Y+1)
                   [\bar Y(\bar Y+1)-1],\\
 F(u)&=E(u)\bigl[(1+2(u^2-1)U^2E(u))b-c\bigr].
\end{align*}
For $j=1,3$, let $Q_j$ denote the ordinary polynomial quotient of $F(u)$ by
$u^j$, evaluated at $u=A$. This is formal polynomial division, never a floor
or rational instruction in the equation. Since
\[
 F(u)=(b-c)E(u)+2U^2b(u^2-1)E(u)^2,
\]
$Q_3$ is coefficientwise even. Exact polynomial algebra gives
\begin{equation}
 \Omega=U(4AEV-2Q_1-Q_3/2).
 \label{eq:weighted-quotient-center}
\end{equation}
The displayed division-free evaluator is precisely this polynomial. For
example, $S$ is the quotient of $E$ by $A$. If
\[
 J_1=S(2E-AS),\quad J_3=e_2(2S-Ae_2),\quad
 H_1=AE^2-J_1,\quad H_3=J_1-J_3,
\]
then $Q_1=(b-c)S+2\theta H_1$, $Q_3=2\theta H_3$, and
\[
 4H_1+H_3=A(4E^2+3S^2+e_2^2)-2S(3E+e_2).
\]
Substitution proves \eqref{eq:weighted-quotient-center} in the free integer
polynomial ring.

For a polynomial in $u$, use a superscript $\#$ for evaluation at $A$ after
taking absolute values of all coefficients. Put
$\delta=(\bar Y+1)/A<1/16$. The inequalities
\[
 |e_1|<2e_2(\bar Y+1),\qquad
 0<e_0\le e_2(\bar Y+1)^2
\]
give
\[
 E^\#(A)\le e_2(A+\bar Y+1)^2,
 \quad E\ge e_2A^2(1-2\delta),\quad E^\#(A)<2E.
\]
Consequently
\[
 F^\#(A)<2E(|b|+c)+8(A^2+1)U^2E^2|b|.
\]
At a zero $|b|\ge1$, $c+1<2U^2E$, and
$|V|>\Delta U^2E|b|$, because $f-c>a=\Delta U^2E$.
It follows that
\[
 \frac{F^\#(A)}{|EV|}
 <\frac4\Delta+\frac{8(A^2+1)}\Delta
 =8+\frac{20}\Delta<10.
\]
Deletion of low coefficients gives $|Q_j|\le F^\#(A)/A^j$. Thus, using
$A\ge16$ in \eqref{eq:weighted-quotient-center},
\[
 |\Omega-4AUEV|
 <10U|EV|\left(\frac2A+\frac1{2A^3}\right)<AU|EV|.
\]
In particular $\Omega$ has the sign of $V$, and the required uniform signed
estimate is
\begin{equation}
 3AUE|V|<|\Omega|<5AUE|V|.
 \label{eq:weighted-center-bound}
\end{equation}

\subsection{The new weighted transport argument}

The source178 weight $b^2$ does not have a proved adequate margin in this free
chart when auxiliary index three is possible. The additional factor $UE$ in
\eqref{eq:weighted-L} resolves that problem. Indeed, from
\eqref{eq:weighted-inherited-bounds} and
\eqref{eq:weighted-center-bound},
\begin{align}
 4gy+|\Omega|
 &<11AUE|V|
 <\frac{33}{2}AUEf|b|\notag\\
 &\le\frac{33A}{8\Delta}\,b^2UE
 <\frac L3.\label{eq:weighted-margin}
\end{align}
The last inequality follows from $99A<8(A^2-1)$ for $A\ge16$; here in fact
$A>1024$. Thus the proof neither assumes nor needs the stronger scaled-chart
index exclusion.

Since $k\ge2$, the inverse-Lucas correction satisfies
\[
 e^*\le k-R-1<b,
 \qquad \rho<\frac{|b|}{q^7}<\frac L2.
\]
This holds for every integer $R$, including zero and negative values.
If $B\le-1$, then $Wt=B/h+LB<-L$. Solving
\eqref{eq:weighted-carrier} for the positive input $n$ and using positive
$\kappa,\sigma$ gives
\[
 n=\varepsilon4gy+\Omega-\kappa-\sigma+Wt+\rho
 <\frac L3-L+\frac L2<0,
\]
a contradiction. Hence $B\ge0$. Since $F>0$ and $K_0+w>0$, the identity
$B=(K_0+w)C-F$ forces $C>0$.

Only now invoke the authenticated whole-tuple mask estimate
$3q<R<q^4-2$. Since $c>4q^4$, it implies $b,V>0$. The unchanged normalized
strong-index and signed auxiliary-index results identify $p=R$, from
\[
 V=-c\pmod f,\qquad V=-R\pmod c,\qquad c^2\mid2Ug.
\]
The retained parity argument yields odd $R$, even $q$, and odd $J$.

\subsection{Authentication, projection, and refutation}

The only changes in the rational carrier are the integer polynomial
$\Omega$ and the integer weight $L$, with $W\equiv1\pmod h$. The parent
authentication and projection proof therefore applies in the following order.
The first Pell norm gives $k=2\psi_{1+2v}(m)$, $m\ge2$, and monotonicity gives
$m<R$. Integral $z$ and the literal denominator give $q^4\mid e^*$.
The full retained Lucas identity is
\[
 3\bigl(e^*-(2m-R-1)\bigr)
   =4v^2(24Z_m-A_mS_m),
\]
where the right side is an integer expression from the finite binomial
expansion of the Pell sequence. Since $v=q^3ws^2$, at every prime dividing
$q$, including3, this implies $q^4\mid(2m-R-1)$. Its absolute value is below
$R<q^4$, so $m=(R+1)/2$.

Reduction of the literal numerator modulo $H$ gives $H\mid6hK$. The
$\lambda=q$ uniform ratio theorem and the complete-product strict bound force
$X=2^R$ and positive integral $\kappa$; this step does not cancel an
unproved factor of $6h$. Because $q\mid X$, the radix $q$ is a power of2.
Rank growth gives $q^6\mid X$, hence $q^8\mid v$. Only now does the same
Lucas identity imply $q^7\mid e^*$, so $\rho$ is integral.

Finally \eqref{eq:weighted-carrier}, with integral $\Omega$ and
$W\equiv1\pmod h$, gives $q^3\mid Y$ and $h\mid B$. Thus $t=B/h$ is a
nonnegative integer. These are exactly the parent native outer conditions.
The inherited fixed-minus parity, half-binomial rounding, and fresh native
reconstruction supply the native words at the complete fixed tuple; its
named syntactic calibration yields a finite ZFC refutation. This does not
replace that calibration by the elementary coefficient inequalities.

\subsection{Constructive positive completeness}

Conversely, a finite syntactic ZFC refutation supplies the parent canonical
outer data and positive $\eta,\zeta,i,T$. In particular $b,V>0$, $t\ge0$ is
integral, and $\rho,\kappa,\sigma$ are integral. The retained canonical
estimates give $g>\kappa+\sigma+|\rho|$. Keep these inputs and define
\[
 n=4gy+\Omega+(1+hL)t+\rho-\kappa-\sigma.
\]
Every term is integral; $\Omega>0$, $y\ge1$, and $t\ge0$ make $n>0$.
Its carrier satisfies $z=4gy$, so \eqref{eq:weighted-root} makes the literal
polynomial vanish. This proves both directions of positive solvability
versus finite refutability. The arithmetic arguments and the retained
proof-code equivalence are internalizable in ZFC; negation gives the
consistency theorem. No restriction to externally standard proof codes is
introduced.

\subsection{Exact degree and finite verification}

The free-chart degree weights for $(A,\bar Y,k,U,b,c)$ are
$(4,2,1,5,4,2)$. Exact sparse collection gives 84 center monomials, degree59,
and 594 monomials in
$D_{\rm norm}=\Omega^2-16GV^2$, of weighted degree 94. Its entire highest
component is
\[
 [D_{\rm norm}]_{94}
 =8A^6\bar Y^8k^8U^6b^2(A^2-6\bar Y^4).
\]
The exact output decomposition is
\[
 P=Q\bigl[Z_0^2-2MZ_0\Omega+M^2D_{\rm norm}\bigr]
      +16M^2G-(Z_0-M\Omega)^2.
\]
Here
\[
 \deg(Q,M,G,V,E,\Omega,L,Z_0)=(72,12,46,36,14,59,27,41).
\]
The five respective degree bounds in the decomposition are
$154,184,190,70,142$, so the norm-defect term uniquely controls the top
degree. With $q_0=\beta J$ and $k_0=\eta+\zeta$, the complete actual
highest component is
\begin{equation}
 [P]_{190}
 =9\cdot2^{11}q_0^{74}w^{18}s^{52}i^{10}k_0^{32}
       (q_0-F)^2(w^2-6s^2).
 \label{eq:positive190-leading}
\end{equation}
This is nonzero for the actual fixed positive $\beta$, proving exact degree 190.
The full leading form has 198 collected monomials.

The independent standard-library checker verifies the actual literal graph,
all-variable/port/gate liveness, the entire center and norm-defect identities,
the complete leading form, all four payload hashes, and 60 deliberately
corrupted center operators. Both normal and optimized Python runs pass.
The separate frozen numeral recipe reconstructs every coefficient byte.
These checks authenticate finite artifacts; the unbounded Pell and logical
theorems remain ordinary mathematical proofs with the stated inherited
dependencies.

\paragraph{Further research.}
The weighted transport argument admits other integer polynomial centers with
the same uniform $O(AUE|V|)$ signed bound. Reducing the 30-gate center
evaluator is a concrete cost target. Reducing the degree requires more than
the next term of the same $A$-Laurent expansion: that term removes the
$8A^8\bar Y^8k^8U^6b^2$ component but retains a nonzero multiple of
$A^6\bar Y^{12}k^8U^6b^2$, of the same weighted degree 94. Any proposed
division by $A$ or $w$ must be replaced by a proved polynomial identity or
an explicitly counted additional witness.


\section{Verification, interpretation, and further research}
\label{sec:verification}

\subsection{What the executable checks establish}

The archive supplies literal polynomial data, pure Python evaluators,
producers, and independent finite checkers. The direct checks collect
every quartic coefficient, verify all declared coordinates occur, bind
the source changes, and replay the exact projection and coordinate
certificates. The compact checks reconstruct the four fixed integers,
check the entire arithmetic graph, collect the free center and norm
defect, and prove the exact degree through a nonzero full leading
homogeneous polynomial. Both ordinary and optimized Python modes are
used; substantive acceptance conditions do not vanish with assertions.
Targeted corruptions exercise the relevant mathematical comparisons.

The main compact center has 84 monomials in its free polynomial ring.
Its norm defect has 594, and weighted degree 94. The resulting
degree-190 leading homogeneous form has 198 monomials when expanded
in the independent highest-degree inputs. These counts are finite
algebraic facts, not probabilistic tests. Auxiliary Pell residue tests,
when retained from a parent, are examples and are not zeros of a
consistency polynomial.

The immediate incoming direct archive is retained byte for byte under
\code{dependencies/Pell_Growth_210_original.zip}. It supplies the original
proofs, source equations, and prior certificates. The new component
directories include the specific inputs they need. Their source ledgers
pin the snapshot and file hashes, and the release manifest authenticates
the delivered bytes. Selected proof-assistant sources are evidence for
the named calibration boundary; they are not represented as a newly
compiled complete dependency closure.

\subsection{Why these are improvements, and what they do not settle}

The new direct equation reduces the number of unrestricted integer
unknowns while retaining degree four. The new compact equation lowers
the exact degree at ten positive unknowns. Its extra operations pay
for a stronger signed transport bound and an exact center; they are
included in the count. Both are concrete tradeoffs against inspected
predecessors. Earlier ProveIt files already contain fewer-variable
representations in other regimes, so the direct count must not be
described as a record over all Diophantine encodings.

The finite-marker theorem is useful beyond this particular count. It
separates finite proof transport from a global bijection of all codes.
It permits the arithmetic of the constructors to be optimized while
unused numbers are handled uniformly as temporary atoms in each finite
proof. This is a rigorously justified design freedom, even when one
particular sparse code yields only a modest numerical improvement.

\subsection{Further research questions}

\begin{enumerate}
\item \textbf{Optimize proof syntax under finite-marker equivalence.}
Can a sparse constructor family create more shared quadratic expressions
than the four offsets in \eqref{eq:sparse-constructors}? An injective
partial coding with strict argument descent is enough; surjectivity
need not be restored. Any new family must still preserve the designated
truth and false codes, the guarded schemata, and the updated growth bounds.
\item \textbf{Change the calculus, not just its numeral coding.}
The largest term is the guarded seed bundle. A shorter calibrated
schema could reduce its depth and the surrounding arithmetic at once.
This requires an actual two-way finite-refutation proof; a smaller
local expression alone is insufficient.
\item \textbf{Couple the transport weight to a sharper center.}
For the free chart, can one lower the weighted norm-defect degree below
94 while keeping both the early signed bound and positive completeness?
Merely removing the index-three case by assumption would invalidate
the present proof. Higher quotient corrections should be compared
with the extra degree they introduce through the carrier.
\item \textbf{Improve the small-coefficient quartic tradeoff.}
The affine slice removes a coordinate but increases coefficient height.
Different primitive affine charts or shared divisibility coordinates
may retain the reduction with a smaller height. A successful change
should report the collected polynomial's height, not only circuit literals.
\item \textbf{Reduce verification dependencies.}
The new sparse recoding is a short finite-tree argument and is a suitable
isolated formalization target if proof-assistant work is resumed.
The weighted transport inequality can also be formalized separately
from the much larger inherited native calibration.
\end{enumerate}

\begin{thebibliography}{9}
\bibitem{repo} Vladimir Reshetnikov, \emph{ProveIt}, snapshot
\code{d123610bc98fef2006bff1c9466bc3f1221003e7},
\url{https://github.com/VladimirReshetnikov/ProveIt/tree/d123610bc98fef2006bff1c9466bc3f1221003e7}.
\bibitem{growth210} \emph{Pell growth from divisibility: an explicit quartic
for Con(ZFC)}, incoming archive \code{ConZFC_Pell_Growth_210_2026-10-25.zip}
at the snapshot above; the complete archive is retained with this paper.
Its standalone growth and affine-radix proofs are the immediate arithmetic
dependencies of the direct construction.
\bibitem{anchor173} \emph{A reused CRT index: an explicit Diophantine
consistency equation}, package 173 in
\code{Computability/DiophantineConsistency/con-zfc-diophantine-forms}.
This supplies the sharp label-product separation and reused-index anchor.
\bibitem{free172} \emph{Pooled powers, round 21}, package 172, free-chart
and Lucas-correction proofs, with the complete pinned native coefficient
tuple. The relevant original text and literal parent are retained under
\code{positive190/inherited}.
\bibitem{center178} \emph{Polynomial centers for Diophantine consistency},
package 178, whole-product center and polynomial quotient identity.
The relevant original proof and algebra are retained under
\code{positive190/inherited}.
\bibitem{signed} \emph{Signed horizons and residue-encoded Pell witnesses},
incoming archive \code{ConZFC_Signed_Horizons_and_Residue_Carriers_2026-10-25.zip}
at the snapshot above. Signed affine horizons are credited to this
predecessor; their composition with the divisibility-growth sign bootstrap
is proved here.
\end{thebibliography}

\appendix
\section{Complete literal residuals of the direct quartic}
\label{app:residuals}

Every coordinate $x_0,\ldots,x_{\DirectLastVariable}$ is an unrestricted
integer. The single equation is
\[
 \sum_{j=0}^{\DirectLastRow}R_j^2=0.
\]
All coefficients below are integers. This list and the machine-readable
residual file define exactly the same polynomial. Subscripts are coordinate
indices and fixed exponents are literal exponents.
% Complete literal residuals; generated from equation.json.
\allowdisplaybreaks[4]
\begin{align*}
% RBEGIN 0
R_{0}={}&2-3x_{3}-4x_{3}x_{108}+x_{6}+4x_{108}\\[3pt]
% REND 0
% RBEGIN 1
R_{1}={}&-x_{1}-x_{3}x_{4}+x_{5}x_{6}\\[3pt]
% REND 1
% RBEGIN 2
R_{2}={}&-1+x_{0}-x_{5}x_{6}\\[3pt]
% REND 2
% RBEGIN 3
R_{3}={}&-2x_{1}x_{135}-x_{2}x_{17}+x_{3}+2x_{15}\\[3pt]
% REND 3
% RBEGIN 4
R_{4}={}&-x_{3}-2x_{3}x_{15}+2x_{15}+x_{20}\\[3pt]
% REND 4
% RBEGIN 5
R_{5}={}&-x_{13}x_{20}+x_{21}\\[3pt]
% REND 5
% RBEGIN 6
R_{6}={}&-2x_{0}^{2}+2x_{0}x_{11}+2x_{0}x_{12}+2x_{0}x_{21}+2x_{0}x_{22}-x_{11}^{2}-2x_{11}x_{21}-x_{12}^{2}\\*
&-2x_{12}x_{22}+x_{16}x_{23}-x_{18}x_{19}-x_{21}^{2}-x_{22}^{2}\\[3pt]
% REND 6
% RBEGIN 7
R_{7}={}&-x_{3}x_{14}+x_{14}-x_{14}x_{20}+x_{22}\\[3pt]
% REND 7
% RBEGIN 8
R_{8}={}&-x_{24}x_{27}+x_{27}+x_{29}\\[3pt]
% REND 8
% RBEGIN 9
R_{9}={}&-x_{2}x_{50}+x_{25}x_{123}-x_{28}-x_{29}x_{50}\\[3pt]
% REND 9
% RBEGIN 10
R_{10}={}&x_{30}x_{123}-x_{31}x_{53}-x_{32}\\[3pt]
% REND 10
% RBEGIN 11
R_{11}={}&x_{33}x_{123}-x_{34}x_{53}-x_{35}\\[3pt]
% REND 11
% RBEGIN 12
R_{12}={}&x_{36}x_{123}-x_{37}x_{53}-x_{38}\\[3pt]
% REND 12
% RBEGIN 13
R_{13}={}&x_{39}x_{123}-x_{40}x_{53}-x_{41}\\[3pt]
% REND 13
% RBEGIN 14
R_{14}={}&x_{42}x_{124}-x_{44}x_{53}-x_{45}\\[3pt]
% REND 14
% RBEGIN 15
R_{15}={}&x_{46}x_{124}-x_{47}x_{53}-x_{48}\\[3pt]
% REND 15
% RBEGIN 16
R_{16}={}&-x_{16}x_{50}+x_{53}\\[3pt]
% REND 16
% RBEGIN 17
R_{17}={}&x_{49}-x_{51}x_{53}-x_{52}\\[3pt]
% REND 17
% RBEGIN 18
R_{18}={}&-x_{53}x_{55}+x_{54}-x_{56}\\[3pt]
% REND 18
% RBEGIN 19
R_{19}={}&-x_{53}x_{58}+x_{57}x_{123}-x_{59}\\[3pt]
% REND 19
% RBEGIN 20
R_{20}={}&x_{60}-2x_{60}x_{62}+x_{63}\\[3pt]
% REND 20
% RBEGIN 21
R_{21}={}&-x_{60}x_{63}-x_{63}+x_{64}\\[3pt]
% REND 21
% RBEGIN 22
R_{22}={}&-1+x_{61}^{2}+x_{63}^{2}-x_{64}^{2}\\[3pt]
% REND 22
% RBEGIN 23
R_{23}={}&1-x_{61}x_{66}-x_{65}+x_{67}\\[3pt]
% REND 23
% RBEGIN 24
R_{24}={}&-1+2x_{65}-x_{65}^{2}+2x_{65}x_{67}+x_{66}^{2}\\[3pt]
% REND 24
% RBEGIN 25
R_{25}={}&1+x_{2}^{2}-2x_{2}x_{61}+x_{68}\\[3pt]
% REND 25
% RBEGIN 26
R_{26}={}&-x_{2}x_{100}+x_{16}+x_{68}x_{69}-x_{99}\\[3pt]
% REND 26
% RBEGIN 27
R_{27}={}&-5-8x_{1}-4x_{1}^{2}-16x_{1}x_{15}+4x_{1}x_{61}-16x_{15}-16x_{15}^{2}+8x_{15}x_{61}+4x_{61}-x_{70}\\[3pt]
% REND 27
% RBEGIN 28
R_{28}={}&-1-2x_{1}x_{100}-4x_{15}x_{100}+x_{24}+x_{70}x_{71}-x_{99}-2x_{100}\\[3pt]
% REND 28
% RBEGIN 29
R_{29}={}&-x_{61}x_{73}-x_{72}+x_{74}\\[3pt]
% REND 29
% RBEGIN 30
R_{30}={}&-1-x_{72}^{2}+2x_{72}x_{74}+x_{73}^{2}\\[3pt]
% REND 30
% RBEGIN 31
R_{31}={}&-x_{61}x_{76}-x_{75}+x_{77}\\[3pt]
% REND 31
% RBEGIN 32
R_{32}={}&-1-x_{75}^{2}+2x_{75}x_{77}+x_{76}^{2}\\[3pt]
% REND 32
% RBEGIN 33
R_{33}={}&-x_{61}x_{79}-x_{78}+x_{80}\\[3pt]
% REND 33
% RBEGIN 34
R_{34}={}&-1-x_{78}^{2}+2x_{78}x_{80}+x_{79}^{2}\\[3pt]
% REND 34
% RBEGIN 35
R_{35}={}&-x_{61}x_{82}-x_{81}+x_{83}\\[3pt]
% REND 35
% RBEGIN 36
R_{36}={}&-1-x_{81}^{2}+2x_{81}x_{83}+x_{82}^{2}\\[3pt]
% REND 36
% RBEGIN 37
R_{37}={}&-x_{61}x_{85}-x_{84}+x_{86}\\[3pt]
% REND 37
% RBEGIN 38
R_{38}={}&-1-x_{84}^{2}+2x_{84}x_{86}+x_{85}^{2}\\[3pt]
% REND 38
% RBEGIN 39
R_{39}={}&-x_{61}x_{88}-x_{87}+x_{89}\\[3pt]
% REND 39
% RBEGIN 40
R_{40}={}&-1-x_{87}^{2}+2x_{87}x_{89}+x_{88}^{2}\\[3pt]
% REND 40
% RBEGIN 41
R_{41}={}&1-x_{61}x_{91}-x_{90}+x_{92}\\[3pt]
% REND 41
% RBEGIN 42
R_{42}={}&-1+2x_{90}-x_{90}^{2}+2x_{90}x_{92}+x_{91}^{2}\\[3pt]
% REND 42
% RBEGIN 43
R_{43}={}&-x_{61}x_{94}-x_{93}+x_{95}\\[3pt]
% REND 43
% RBEGIN 44
R_{44}={}&-1-x_{93}^{2}+2x_{93}x_{95}+x_{94}^{2}\\[3pt]
% REND 44
% RBEGIN 45
R_{45}={}&-x_{61}x_{97}-x_{96}+x_{98}\\[3pt]
% REND 45
% RBEGIN 46
R_{46}={}&-1-x_{96}^{2}+2x_{96}x_{98}+x_{97}^{2}\\[3pt]
% REND 46
% RBEGIN 47
R_{47}={}&1-x_{61}x_{100}-x_{99}+x_{101}\\[3pt]
% REND 47
% RBEGIN 48
R_{48}={}&-1+2x_{99}-x_{99}^{2}+2x_{99}x_{101}+x_{100}^{2}\\[3pt]
% REND 48
% RBEGIN 49
R_{49}={}&-x_{61}x_{103}-x_{102}+x_{104}\\[3pt]
% REND 49
% RBEGIN 50
R_{50}={}&-1-x_{102}^{2}+2x_{102}x_{104}+x_{103}^{2}\\[3pt]
% REND 50
% RBEGIN 51
R_{51}={}&-x_{61}x_{106}-x_{105}+x_{107}\\[3pt]
% REND 51
% RBEGIN 52
R_{52}={}&-1-x_{105}^{2}+2x_{105}x_{107}+x_{106}^{2}\\[3pt]
% REND 52
% RBEGIN 53
R_{53}={}&-28-x_{0}^{2}-12x_{1}-6x_{1}^{2}-16x_{1}x_{15}-8x_{1}x_{108}-x_{2}^{2}-x_{3}^{2}-x_{7}^{2}-x_{8}^{2}-x_{9}^{2}\\*
&-x_{10}^{2}-x_{11}^{2}-x_{12}^{2}-x_{13}^{2}-x_{14}^{2}-16x_{15}-20x_{15}^{2}-x_{16}^{2}+2x_{24}-2x_{24}^{2}\\*
&-x_{25}^{2}-x_{26}^{2}-x_{30}^{2}-x_{33}^{2}-x_{36}^{2}-x_{39}^{2}-x_{42}^{2}-x_{43}^{2}-x_{46}^{2}-x_{49}^{2}\\*
&-x_{50}^{2}-x_{54}^{2}-x_{57}^{2}+x_{60}-72x_{108}-64x_{108}^{2}\\[3pt]
% REND 53
% RBEGIN 54
R_{54}={}&-x_{61}x_{110}+x_{119}\\[3pt]
% REND 54
% RBEGIN 55
R_{55}={}&-1+x_{109}^{2}+x_{110}^{2}-x_{119}^{2}\\[3pt]
% REND 55
% RBEGIN 56
R_{56}={}&-x_{112}x_{113}+x_{120}\\[3pt]
% REND 56
% RBEGIN 57
R_{57}={}&-1+x_{111}^{2}+x_{112}^{2}-x_{120}^{2}\\[3pt]
% REND 57
% RBEGIN 58
R_{58}={}&-x_{106}^{2}+x_{121}\\[3pt]
% REND 58
% RBEGIN 59
R_{59}={}&x_{110}-2x_{115}x_{121}-x_{121}\\[3pt]
% REND 59
% RBEGIN 60
R_{60}={}&-1-4x_{110}-8x_{110}x_{114}+x_{113}\\[3pt]
% REND 60
% RBEGIN 61
R_{61}={}&-x_{61}-x_{109}x_{116}+x_{113}\\[3pt]
% REND 61
% RBEGIN 62
R_{62}={}&-x_{107}-x_{109}x_{117}+x_{111}\\[3pt]
% REND 62
% RBEGIN 63
R_{63}={}&-2x_{15}-4x_{106}x_{118}+x_{112}\\[3pt]
% REND 63
% RBEGIN 64
R_{64}={}&17+x_{0}^{2}-2x_{0}x_{100}+4x_{1}+x_{1}^{2}-2x_{1}x_{73}+8x_{1}x_{108}+x_{3}^{2}-2x_{3}x_{91}+x_{7}^{2}\\*
&-2x_{7}x_{76}+x_{8}^{2}-2x_{8}x_{79}+x_{9}^{2}-2x_{9}x_{82}+x_{10}^{2}-2x_{10}x_{85}+x_{11}^{2}-2x_{11}x_{88}\\*
&+x_{12}^{2}-2x_{12}x_{94}+x_{13}^{2}-2x_{13}x_{97}+x_{14}^{2}-2x_{14}x_{103}-4x_{15}x_{106}-x_{61}x_{122}\\*
&-4x_{66}-8x_{66}x_{108}+x_{67}^{2}-4x_{73}-8x_{73}x_{108}-2x_{74}+x_{74}^{2}-2x_{77}+x_{77}^{2}-2x_{80}\\*
&+x_{80}^{2}-2x_{83}+x_{83}^{2}-2x_{86}+x_{86}^{2}-2x_{89}+x_{89}^{2}+x_{92}^{2}-2x_{95}+x_{95}^{2}-2x_{98}\\*
&+x_{98}^{2}+x_{101}^{2}-2x_{104}+x_{104}^{2}+2x_{106}^{2}-2x_{107}+x_{107}^{2}+32x_{108}+32x_{108}^{2}+x_{122}\\*
&-x_{125}^{2}-x_{126}^{2}-x_{127}-x_{127}^{2}\\[3pt]
% REND 64
% RBEGIN 65
R_{65}={}&-x_{26}x_{123}+x_{50}\\[3pt]
% REND 65
% RBEGIN 66
R_{66}={}&-x_{43}x_{124}+x_{50}\\[3pt]
% REND 66
% RBEGIN 67
R_{67}={}&-1-4x_{15}^{2}-x_{66}^{2}-x_{73}^{2}-x_{76}^{2}-x_{79}^{2}-x_{82}^{2}-x_{85}^{2}-x_{88}^{2}-x_{91}^{2}\\*
&-x_{94}^{2}-x_{97}^{2}-x_{100}^{2}-x_{103}^{2}+x_{106}^{2}-x_{125}^{2}-x_{126}^{2}-x_{127}-x_{127}^{2}\\[3pt]
% REND 67
% RBEGIN 68
R_{68}={}&-3-4x_{28}^{2}-4x_{32}^{2}-4x_{35}^{2}-4x_{38}^{2}-4x_{41}^{2}-4x_{45}^{2}-4x_{48}^{2}+x_{50}^{2}-4x_{52}^{2}\\*
&-4x_{56}^{2}-4x_{59}^{2}-x_{128}^{2}-x_{129}^{2}+3x_{130}-x_{130}^{2}\\[3pt]
% REND 68
% RBEGIN 69
R_{69}={}&x_{24}x_{132}-2x_{25}x_{73}-2x_{30}x_{76}-2x_{33}x_{79}-2x_{36}x_{82}-2x_{39}x_{85}-2x_{42}x_{88}\\*
&-2x_{46}x_{94}-2x_{49}x_{97}-2x_{54}x_{103}-2x_{57}x_{106}-2x_{61}x_{132}+2x_{73}x_{74}+2x_{76}x_{77}\\*
&+2x_{79}x_{80}+2x_{82}x_{83}+2x_{85}x_{86}+2x_{88}x_{89}+2x_{94}x_{95}+2x_{97}x_{98}+2x_{103}x_{104}\\*
&+2x_{106}x_{107}+x_{131}\\[3pt]
% REND 69
% RBEGIN 70
R_{70}={}&-10-x_{24}x_{131}+x_{25}^{2}-2x_{25}x_{72}+x_{30}^{2}-2x_{30}x_{75}+x_{33}^{2}-2x_{33}x_{78}+x_{36}^{2}\\*
&-2x_{36}x_{81}+x_{39}^{2}-2x_{39}x_{84}+x_{42}^{2}-2x_{42}x_{87}+x_{46}^{2}-2x_{46}x_{93}+x_{49}^{2}\\*
&-2x_{49}x_{96}+x_{54}^{2}-2x_{54}x_{102}+x_{57}^{2}-2x_{57}x_{105}+2x_{72}x_{74}+2x_{75}x_{77}+2x_{78}x_{80}\\*
&+2x_{81}x_{83}+2x_{84}x_{86}+2x_{87}x_{89}+2x_{93}x_{95}+2x_{96}x_{98}+2x_{102}x_{104}+2x_{105}x_{107}+x_{132}\\[3pt]
% REND 70
% RBEGIN 71
R_{71}={}&x_{24}x_{134}-2x_{26}x_{66}-2x_{43}x_{91}-2x_{50}x_{100}-2x_{61}x_{134}+2x_{66}+2x_{66}x_{67}+2x_{91}\\*
&+2x_{91}x_{92}+2x_{100}+2x_{100}x_{101}+x_{133}-x_{134}\\[3pt]
% REND 71
% RBEGIN 72
R_{72}={}&-2-x_{24}x_{133}+x_{26}^{2}-2x_{26}x_{65}+x_{43}^{2}-2x_{43}x_{90}+x_{50}^{2}-2x_{50}x_{99}+2x_{65}\\*
&+2x_{65}x_{67}+2x_{90}+2x_{90}x_{92}+x_{99}^{2}-x_{100}^{2}+x_{133}+x_{134}\\[3pt]
% REND 72
% RBEGIN 73
R_{73}={}&-1-x_{1}-2x_{1}x_{135}+x_{3}\\[3pt]
% REND 73
% RBEGIN 74
R_{74}={}&-2x_{1}-4x_{1}x_{108}+x_{137}\\[3pt]
% REND 74
% RBEGIN 75
R_{75}={}&-1-16x_{1}x_{137}+x_{136}^{2}-64x_{137}^{2}\\[3pt]
% REND 75
% RBEGIN 76
R_{76}={}&-1048576+x_{1}-2097152x_{108}-4194304x_{108}x_{138}-2097152x_{138}\\[3pt]
% REND 76
% RBEGIN 77
R_{77}={}&-x_{7}^{2}-2x_{7}x_{8}-x_{8}^{2}+x_{139}\\[3pt]
% REND 77
% RBEGIN 78
R_{78}={}&-16-32x_{7}-16x_{7}^{2}-32x_{7}x_{139}-32x_{139}-16x_{139}^{2}+x_{140}\\[3pt]
% REND 78
% RBEGIN 79
R_{79}={}&-144-408x_{7}-289x_{7}^{2}-544x_{7}x_{139}-136x_{7}x_{140}-384x_{139}-256x_{139}^{2}-128x_{139}x_{140}\\*
&-96x_{140}-16x_{140}^{2}+x_{141}\\[3pt]
% REND 79
% RBEGIN 80
R_{80}={}&-9-24x_{7}-16x_{7}^{2}-32x_{7}x_{141}-24x_{141}-16x_{141}^{2}+x_{142}\\[3pt]
% REND 80
% RBEGIN 81
R_{81}={}&-x_{8}^{2}-2x_{8}x_{9}-x_{9}^{2}+x_{143}\\[3pt]
% REND 81
% RBEGIN 82
R_{82}={}&-x_{7}^{2}-2x_{7}x_{9}-x_{9}^{2}+x_{144}\\[3pt]
% REND 82
% RBEGIN 83
R_{83}={}&-16-64x_{7}-65x_{7}^{2}-8x_{7}x_{8}-64x_{7}x_{139}-8x_{7}x_{143}-64x_{7}x_{144}-16x_{8}^{2}-32x_{8}x_{143}\\*
&-32x_{139}-16x_{139}^{2}-32x_{139}x_{144}-16x_{143}^{2}-32x_{144}-16x_{144}^{2}+x_{145}\\[3pt]
% REND 83
% RBEGIN 84
R_{84}={}&-100-400x_{7}-404x_{7}^{2}-32x_{7}x_{8}-640x_{7}x_{139}-32x_{7}x_{143}-160x_{7}x_{145}-64x_{8}^{2}\\*
&-128x_{8}x_{143}-320x_{139}-256x_{139}^{2}-128x_{139}x_{145}-64x_{143}^{2}-80x_{145}-16x_{145}^{2}+x_{146}\\[3pt]
% REND 84
% RBEGIN 85
R_{85}={}&-4-32x_{9}-64x_{9}^{2}-64x_{9}x_{143}-64x_{9}x_{144}-16x_{143}-16x_{143}^{2}-32x_{143}x_{144}-16x_{144}\\*
&-16x_{144}^{2}+x_{147}\\[3pt]
% REND 85
% RBEGIN 86
R_{86}={}&-49-56x_{7}-16x_{7}^{2}-128x_{7}x_{9}-32x_{7}x_{139}-128x_{7}x_{144}-32x_{7}x_{147}-224x_{9}-256x_{9}^{2}\\*
&-128x_{9}x_{139}-512x_{9}x_{144}-128x_{9}x_{147}-56x_{139}-16x_{139}^{2}-128x_{139}x_{144}-32x_{139}x_{147}\\*
&-224x_{144}-256x_{144}^{2}-128x_{144}x_{147}-56x_{147}-16x_{147}^{2}+x_{148}\\[3pt]
% REND 86
% RBEGIN 87
R_{87}={}&-4-16x_{7}-17x_{7}^{2}-16x_{7}x_{8}-8x_{7}x_{143}-32x_{7}x_{144}-4x_{8}-17x_{8}^{2}-32x_{8}x_{143}\\*
&-8x_{8}x_{144}-16x_{143}^{2}-16x_{144}-16x_{144}^{2}+x_{149}\\[3pt]
% REND 87
% RBEGIN 88
R_{88}={}&-400-800x_{7}-404x_{7}^{2}-2752x_{7}x_{8}-2592x_{7}x_{143}-160x_{7}x_{149}-2720x_{8}-4688x_{8}^{2}\\*
&-8832x_{8}x_{143}-544x_{8}x_{149}-2560x_{143}-4160x_{143}^{2}-512x_{143}x_{149}-160x_{149}-16x_{149}^{2}\\*
&+x_{150}\\[3pt]
% REND 88
% RBEGIN 89
R_{89}={}&-4-16x_{7}-16x_{7}^{2}-32x_{7}x_{9}-32x_{7}x_{139}-32x_{7}x_{143}-16x_{9}-16x_{9}^{2}-32x_{9}x_{139}\\*
&-32x_{9}x_{143}-16x_{139}-16x_{139}^{2}-32x_{139}x_{143}-16x_{143}-16x_{143}^{2}+x_{151}\\[3pt]
% REND 89
% RBEGIN 90
R_{90}={}&-49-280x_{9}-400x_{9}^{2}-640x_{9}x_{143}-160x_{9}x_{144}-160x_{9}x_{151}-224x_{143}-256x_{143}^{2}\\*
&-128x_{143}x_{144}-128x_{143}x_{151}-56x_{144}-16x_{144}^{2}-32x_{144}x_{151}-56x_{151}-16x_{151}^{2}+x_{152}\\[3pt]
% REND 90
% RBEGIN 91
R_{91}={}&-400-600x_{7}-225x_{7}^{2}-480x_{7}x_{139}+120x_{7}x_{140}-640x_{139}-256x_{139}^{2}+128x_{139}x_{140}\\*
&+160x_{140}-16x_{140}^{2}+x_{153}\\[3pt]
% REND 91
% RBEGIN 92
R_{92}={}&-9-24x_{7}-16x_{7}^{2}-32x_{7}x_{153}-24x_{153}-16x_{153}^{2}+x_{154}\\[3pt]
% REND 92
% RBEGIN 93
R_{93}={}&-64-256x_{7}-256x_{7}^{2}-32x_{7}x_{8}-512x_{7}x_{153}-128x_{7}x_{154}-16x_{8}-x_{8}^{2}-32x_{8}x_{153}\\*
&-8x_{8}x_{154}-256x_{153}-256x_{153}^{2}-128x_{153}x_{154}-64x_{154}-16x_{154}^{2}+x_{155}\\[3pt]
% REND 93
% RBEGIN 94
R_{94}={}&-33489-87840x_{7}-57600x_{7}^{2}+1920x_{7}x_{9}-115200x_{7}x_{139}+7680x_{7}x_{140}+1920x_{7}x_{143}+1464x_{9}\\*
&-16x_{9}^{2}+1920x_{9}x_{139}-128x_{9}x_{140}-32x_{9}x_{143}-87840x_{139}-57600x_{139}^{2}+7680x_{139}x_{140}\\*
&+1920x_{139}x_{143}+5856x_{140}-256x_{140}^{2}-128x_{140}x_{143}+1464x_{143}-16x_{143}^{2}+x_{156}\\[3pt]
% REND 94
% RBEGIN 95
R_{95}={}&-225-600x_{9}-400x_{9}^{2}-640x_{9}x_{143}-160x_{9}x_{144}-160x_{9}x_{156}-480x_{143}-256x_{143}^{2}\\*
&-128x_{143}x_{144}-128x_{143}x_{156}-120x_{144}-16x_{144}^{2}-32x_{144}x_{156}-120x_{156}-16x_{156}^{2}\\*
&+x_{157}\\[3pt]
% REND 95
% RBEGIN 96
R_{96}={}&-x_{9}^{2}-2x_{9}x_{10}-x_{10}^{2}+x_{158}\\[3pt]
% REND 96
% RBEGIN 97
R_{97}={}&-17x_{7}^{2}-32x_{7}x_{8}-8x_{7}x_{9}-32x_{7}x_{143}-32x_{7}x_{144}-8x_{7}x_{158}-16x_{8}^{2}-32x_{8}x_{143}\\*
&-32x_{8}x_{144}-16x_{9}^{2}-32x_{9}x_{158}-16x_{143}^{2}-32x_{143}x_{144}-16x_{144}^{2}-16x_{158}^{2}+x_{159}\\[3pt]
% REND 97
% RBEGIN 98
R_{98}={}&-260x_{7}^{2}-544x_{7}x_{9}-512x_{7}x_{144}-544x_{7}x_{158}-128x_{7}x_{159}-320x_{9}^{2}-512x_{9}x_{144}\\*
&-640x_{9}x_{158}-128x_{9}x_{159}-256x_{144}^{2}-512x_{144}x_{158}-128x_{144}x_{159}-320x_{158}^{2}\\*
&-128x_{158}x_{159}-16x_{159}^{2}+x_{160}\\[3pt]
% REND 98
% RBEGIN 99
R_{99}={}&-16-136x_{9}-289x_{9}^{2}-544x_{9}x_{144}-136x_{9}x_{147}-128x_{144}-256x_{144}^{2}-128x_{144}x_{147}\\*
&-32x_{147}-16x_{147}^{2}+x_{161}\\[3pt]
% REND 99
% RBEGIN 100
R_{100}={}&-4-16x_{8}-16x_{8}^{2}-32x_{8}x_{143}-16x_{143}-16x_{143}^{2}+x_{162}\\[3pt]
% REND 100
% RBEGIN 101
R_{101}={}&-16-128x_{8}-256x_{8}^{2}-32x_{8}x_{9}-512x_{8}x_{143}-128x_{8}x_{162}-8x_{9}-x_{9}^{2}-32x_{9}x_{143}\\*
&-8x_{9}x_{162}-128x_{143}-256x_{143}^{2}-128x_{143}x_{162}-32x_{162}-16x_{162}^{2}+x_{163}\\[3pt]
% REND 101
% RBEGIN 102
R_{102}={}&-x_{8}^{2}-8x_{8}x_{9}-8x_{8}x_{163}-16x_{9}^{2}-32x_{9}x_{163}-16x_{163}^{2}+x_{164}\\[3pt]
% REND 102
% RBEGIN 103
R_{103}={}&-9-24x_{8}-16x_{8}^{2}-32x_{8}x_{164}-24x_{164}-16x_{164}^{2}+x_{165}\\[3pt]
% REND 103
% RBEGIN 104
R_{104}={}&-64-256x_{8}-256x_{8}^{2}-32x_{8}x_{9}-512x_{8}x_{164}-128x_{8}x_{165}-16x_{9}-x_{9}^{2}-32x_{9}x_{164}\\*
&-8x_{9}x_{165}-256x_{164}-256x_{164}^{2}-128x_{164}x_{165}-64x_{165}-16x_{165}^{2}+x_{166}\\[3pt]
% REND 104
% RBEGIN 105
R_{105}={}&-x_{8}^{2}-8x_{8}x_{9}-8x_{8}x_{166}-16x_{9}^{2}-32x_{9}x_{166}-16x_{166}^{2}+x_{167}\\[3pt]
% REND 105
% RBEGIN 106
R_{106}={}&-16-32x_{8}-16x_{8}^{2}-32x_{8}x_{9}-32x_{8}x_{161}-32x_{8}x_{167}-32x_{9}-16x_{9}^{2}-32x_{9}x_{161}\\*
&-32x_{9}x_{167}-32x_{161}-16x_{161}^{2}-32x_{161}x_{167}-32x_{167}-16x_{167}^{2}+x_{168}\\[3pt]
% REND 106
% RBEGIN 107
R_{107}={}&-64-16x_{7}-x_{7}^{2}-32x_{7}x_{9}-32x_{7}x_{161}-8x_{7}x_{168}-256x_{9}-256x_{9}^{2}-512x_{9}x_{161}\\*
&-128x_{9}x_{168}-256x_{161}-256x_{161}^{2}-128x_{161}x_{168}-64x_{168}-16x_{168}^{2}+x_{169}\\[3pt]
% REND 107
% RBEGIN 108
R_{108}={}&-9-24x_{7}-16x_{7}^{2}-32x_{7}x_{169}-24x_{169}-16x_{169}^{2}+x_{170}\\[3pt]
% REND 108
% RBEGIN 109
R_{109}={}&-64-256x_{7}-256x_{7}^{2}-32x_{7}x_{9}-512x_{7}x_{169}-128x_{7}x_{170}-16x_{9}-x_{9}^{2}-32x_{9}x_{169}\\*
&-8x_{9}x_{170}-256x_{169}-256x_{169}^{2}-128x_{169}x_{170}-64x_{170}-16x_{170}^{2}+x_{171}\\[3pt]
% REND 109
% RBEGIN 110
R_{110}={}&-x_{7}^{2}-8x_{7}x_{9}-8x_{7}x_{171}-16x_{9}^{2}-32x_{9}x_{171}-16x_{171}^{2}+x_{172}\\[3pt]
% REND 110
% RBEGIN 111
R_{111}={}&-9-24x_{7}-16x_{7}^{2}-32x_{7}x_{172}-24x_{172}-16x_{172}^{2}+x_{173}\\[3pt]
% REND 111
% RBEGIN 112
R_{112}={}&-81-288x_{7}-256x_{7}^{2}-128x_{7}x_{8}-128x_{7}x_{143}-512x_{7}x_{172}-128x_{7}x_{173}-72x_{8}-16x_{8}^{2}\\*
&-32x_{8}x_{143}-128x_{8}x_{172}-32x_{8}x_{173}-72x_{143}-16x_{143}^{2}-128x_{143}x_{172}-32x_{143}x_{173}\\*
&-288x_{172}-256x_{172}^{2}-128x_{172}x_{173}-72x_{173}-16x_{173}^{2}+x_{174}\\[3pt]
% REND 112
% RBEGIN 113
R_{113}={}&-25-160x_{8}-305x_{8}^{2}-544x_{8}x_{143}-32x_{8}x_{167}-136x_{8}x_{174}-128x_{143}-256x_{143}^{2}\\*
&-128x_{143}x_{174}-24x_{167}-16x_{167}^{2}-32x_{174}-16x_{174}^{2}+x_{175}\\[3pt]
% REND 113
% RBEGIN 114
R_{114}={}&-64-304x_{8}-336x_{8}^{2}-512x_{8}x_{167}-160x_{8}x_{175}-224x_{167}-192x_{167}^{2}-128x_{167}x_{175}\\*
&-84x_{175}-16x_{175}^{2}+x_{176}\\[3pt]
% REND 114
% RBEGIN 115
R_{115}={}&-73-24x_{7}-16x_{7}^{2}-32x_{7}x_{9}-32x_{7}x_{139}-32x_{7}x_{178}-256x_{8}-256x_{8}^{2}-32x_{8}x_{9}\\*
&-128x_{8}x_{176}-40x_{9}-17x_{9}^{2}-32x_{9}x_{139}-8x_{9}x_{176}-32x_{9}x_{178}-24x_{139}-16x_{139}^{2}\\*
&-32x_{139}x_{178}-64x_{176}-16x_{176}^{2}+x_{177}-24x_{178}-16x_{178}^{2}\\[3pt]
% REND 115
% RBEGIN 116
R_{116}={}&-1-8x_{7}-16x_{7}^{2}-8x_{7}x_{9}-32x_{7}x_{139}-2x_{9}-x_{9}^{2}-8x_{9}x_{139}-8x_{139}-16x_{139}^{2}+x_{178}\\[3pt]
% REND 116
% RBEGIN 117
R_{117}={}&-72-288x_{7}-320x_{7}^{2}-256x_{7}x_{9}-640x_{7}x_{139}-128x_{7}x_{177}-128x_{7}x_{178}-144x_{9}-80x_{9}^{2}\\*
&-256x_{9}x_{139}-32x_{9}x_{177}-128x_{9}x_{178}-288x_{139}-320x_{139}^{2}-128x_{139}x_{177}-128x_{139}x_{178}\\*
&-48x_{177}-16x_{177}^{2}-96x_{178}-64x_{178}^{2}+x_{179}\\[3pt]
% REND 117
% RBEGIN 118
R_{118}={}&-289-2176x_{7}-4096x_{7}^{2}-8192x_{7}x_{139}-512x_{7}x_{179}-2176x_{139}-4096x_{139}^{2}-512x_{139}x_{179}\\*
&-136x_{179}-16x_{179}^{2}+x_{180}\\[3pt]
% REND 118
% RBEGIN 119
R_{119}={}&-x_{1}x_{12}+x_{181}\\[3pt]
% REND 119
% RBEGIN 120
R_{120}={}&-x_{1}x_{11}+x_{182}\\[3pt]
% REND 120
% RBEGIN 121
R_{121}={}&16x_{7}x_{181}-x_{11}+x_{12}+16x_{141}x_{181}+4x_{142}x_{181}+8x_{181}+x_{183}\\[3pt]
% REND 121
% RBEGIN 122
R_{122}={}&16x_{7}x_{181}-x_{11}+x_{12}+4x_{146}x_{181}+8x_{181}+x_{184}\\[3pt]
% REND 122
% RBEGIN 123
R_{123}={}&16x_{7}x_{181}-x_{11}+x_{12}+16x_{139}x_{181}+4x_{148}x_{181}+12x_{181}+x_{185}\\[3pt]
% REND 123
% RBEGIN 124
R_{124}={}&80x_{7}x_{181}+256x_{8}x_{181}-x_{11}+x_{12}+256x_{143}x_{181}+4x_{150}x_{181}+72x_{181}+x_{186}\\[3pt]
% REND 124
% RBEGIN 125
R_{125}={}&16x_{9}x_{181}-x_{11}+x_{12}+16x_{144}x_{181}+4x_{152}x_{181}+12x_{181}+x_{187}\\[3pt]
% REND 125
% RBEGIN 126
R_{126}={}&64x_{7}x_{181}-x_{11}+x_{12}+64x_{153}x_{181}+16x_{154}x_{181}+4x_{155}x_{181}+32x_{181}+x_{188}\\[3pt]
% REND 126
% RBEGIN 127
R_{127}={}&16x_{9}x_{181}-x_{11}+x_{12}+16x_{144}x_{181}+4x_{157}x_{181}+12x_{181}+x_{189}\\[3pt]
% REND 127
% RBEGIN 128
R_{128}={}&64x_{9}x_{181}-x_{11}+x_{12}+64x_{158}x_{181}+4x_{160}x_{181}+x_{190}\\[3pt]
% REND 128
% RBEGIN 129
R_{129}={}&-x_{7}x_{8}+x_{7}x_{9}+x_{8}x_{9}-x_{9}^{2}+x_{191}\\[3pt]
% REND 129
% RBEGIN 130
R_{130}={}&1+2x_{10}+x_{10}^{2}-x_{191}^{2}+x_{192}\\[3pt]
% REND 130
% RBEGIN 131
R_{131}={}&256x_{7}x_{181}-x_{11}+x_{12}+256x_{139}x_{181}+16x_{179}x_{181}+4x_{180}x_{181}+64x_{181}+x_{193}\\[3pt]
% REND 131
% RBEGIN 132
R_{132}={}&x_{18}-x_{192}^{2}-x_{193}^{2}\\[3pt]
% REND 132
% RBEGIN 133
R_{133}={}&-1-4x_{1}x_{7}-4x_{1}x_{139}+x_{194}\\[3pt]
% REND 133
% RBEGIN 134
R_{134}={}&-x_{7}x_{182}-x_{11}+x_{195}\\[3pt]
% REND 134
% RBEGIN 135
R_{135}={}&-x_{8}x_{181}-x_{12}+x_{194}x_{195}+x_{196}\\[3pt]
% REND 135
% RBEGIN 136
R_{136}={}&4x_{7}x_{181}-x_{8}x_{182}-x_{11}+x_{12}+4x_{139}x_{181}+2x_{181}+x_{197}\\[3pt]
% REND 136
% RBEGIN 137
R_{137}={}&-x_{11}x_{183}+x_{12}x_{183}+x_{198}\\[3pt]
% REND 137
% RBEGIN 138
R_{138}={}&-x_{184}x_{198}+x_{199}\\[3pt]
% REND 138
% RBEGIN 139
R_{139}={}&-x_{185}x_{199}+x_{200}\\[3pt]
% REND 139
% RBEGIN 140
R_{140}={}&-x_{186}x_{200}+x_{201}\\[3pt]
% REND 140
% RBEGIN 141
R_{141}={}&-x_{187}x_{201}+x_{202}\\[3pt]
% REND 141
% RBEGIN 142
R_{142}={}&-x_{188}x_{202}+x_{203}\\[3pt]
% REND 142
% RBEGIN 143
R_{143}={}&-x_{189}x_{203}+x_{204}\\[3pt]
% REND 143
% RBEGIN 144
R_{144}={}&-x_{190}x_{204}+x_{205}\\[3pt]
% REND 144
% RBEGIN 145
R_{145}={}&-x_{196}x_{205}+x_{206}\\[3pt]
% REND 145
% RBEGIN 146
R_{146}={}&-x_{197}x_{206}+x_{207}\\[3pt]
% REND 146
% RBEGIN 147
R_{147}={}&x_{19}-x_{207}^{2}\\[3pt]
% REND 147
\end{align*}


\section{Complete fixed coefficients for the ten-variable polynomial}
\label{app:coefficients}
\subsection{The complete fixed coefficient recipe}
All quantities in this subsection are part of a finite specification of fixed
integers; none is an additional unknown in the polynomial.
The symbol $K$ in this inherited recipe denotes
our fixed port $K_0$; the temporary $B$ below is local to this recipe. The
calibrated tuple is copied as a whole.\par
\begingroup
\providecommand{\fl}[1]{\left\lfloor#1\right\rfloor}
\providecommand{\rg}[2]{\{#1,\ldots,#2\}}
\textbf{The four fixed integers.}
Every index, set, sum, product and floor below defines a constant, not
another unknown. Put
\[
\rho=2^{67},\quad B=\rho^{8807},\quad\beta=B-1,\quad
G_m(x)=\sum_{j=0}^{m-1}x^j.
\]
Define
\[
\begin{split}
\mathcal E={}&\rg0{942}\cup\rg{983}{992}\cup\rg{1003}{1011}\\
&\cup\rg{1013}{1031}\cup\{1063,1064,1095,2195,3296\}.
\end{split}
\]
Then
\begin{align*}
M_C={}&B-1-\sum_{e\in\mathcal E\setminus\{1\}}\rho^e-2\rho^{1063},\\
M_F={}&B-1+4+(2^{51}-2)(\rho^{4393}+\rho^{5495})\\
&+\rho^{1013}G_{10}(\rho)+\rho^{1065}G_9(\rho)
  +\rho^{2195}+\rho^{3296}\\
&+(2^{65}-1)\rho^{8792}G_{13}(\rho)+(2^{19}-1)\rho^{8805},\\
K={}&\rho^{30}+\rho^{62}+\rho^{1100}+\rho^{2201}+\rho^{8849}\\
&+(\rho^{4393}+\rho^{5495})\Phi(1/\rho),
\end{align*}
where
\[
\begin{split}
\Phi(x)={}&(1+2^{51})G_{943}(x)
 -2^{10}\sum_{i=0}^{942}v(t_i)x^i\\
&+\sum_{(b,e,m)\in\mathcal C}2^b x^e G_m(2x),\\
\mathcal C={}&\{(10,1023,8),(18,1013,9),(27,1003,8),\\
&(35,983,9),(44,1031,1),(45,1022,1),(46,1011,1),\\
&(47,992,1),(48,1095,1),(49,2195,1),(50,3296,1)\}.
\end{split}
\]
The rational intermediate in K cancels exactly; all four results are positive integers.

\textbf{The 943 indexed fixed tiles.}
Use the actions $u_j,v_j,d_j$ decoded below. Let
\[
\begin{gathered}
h(q,b)=2(q+1)+b,\quad
\mathcal R=\{u_j:d_j=1,\ 0\le j<430\},\\
\mathcal L=\{u_j:d_j=0,\ 0\le j<430\},\\
r_q=6+|\mathcal R\cap\rg0{q-1}|,\quad
l_q=205+|\mathcal L\cap\rg0{q-1}|.
\end{gathered}
\]
An interval ending at $-1$ is empty. The table defines $t_i$ for
$0\le i\le942$; all index ranges are fixed:
\[
\setlength{\arraycolsep}{1.4pt}\begin{array}{c|l}
\text{index }i&\text{tile }t_i=(w,n,e,s)\\\hline
0,1&(2,435,3,1),\ (5,0,0,0)\\
j+2\ (j=2q+b<430)&((1-d_j)l_{u_j},h(q,b),d_jr_{u_j},v_j)\\
432+b\ (b=0,1)&(0,h(215,b),0,h(215,b))\\
434+b\ (b=0,1)&(0,b,0,b)\\
424+2r_q+b\ (q\in\mathcal R,b=0,1)&(r_q,b,0,h(q,b))\\
424+2l_q+b\ (q\in\mathcal L,b=0,1)&(0,b,l_q,h(q,b))\\
924+j\ (0\le j<4)&(0,j+430\fl{j/2},0,438)\\
928&(2,434,2,0)\\
929&(3,436,4,4)\\
930&(4,436,4,0)\\
497+t\ (t=434,435,436)&(1,438,1,t)\\
\substack{934+3a+b-6\fl{a/4}\\a\in\{0,1,4\},\ b=0,1,2}&(a,437,b,437)
\end{array}
\]
For $t=(w,n,e,s)$ put
\[
v(t)=w+2^8n+2^{17}e+2^{25}s+127\,2^{34}.
\]

\textbf{A closed decoder for the 430 actions.}
Put $D_b(A,k)=\fl{A/b^k}\bmod b$, where
$a\bmod b=a-b\fl{a/b}$. For $0\le q<215$, $b=0,1$, $j=2q+b$:
\[
\begin{gathered}
N=\mathcal A\bmod2^{636},\quad
T_0=\fl{\mathcal A/2^{636}}\bmod2^{1675},\\
E_0=\fl{\mathcal A/2^{2311}},\\
\nu_j=D_2(N,j),\quad \sigma_j=\sum_{k=0}^{j-1}\nu_k,\\
m_q=D_2(N,430+q),\quad \tau_q=\sum_{k=0}^{q-1}m_k,\\
F_q=14(1-m_q)+m_qD_{16}(E_0,\tau_q),\\
u_j=\nu_j(1+\sigma_j)+(1-\nu_j)D_{216}(T_0,j-\sigma_j),\\
v_j=D_2(F_q,2b),\quad d_j=D_2(F_q,2b+1).
\end{gathered}
\]
Empty sums are zero. Here $(v_0,d_0,u_0)=(1,1,1)$; the constructor tile
uses this actual first target, including its state renaming.

\textbf{The integer $\mathcal A$.}
The alphabet, in order of digit values $0,\ldots,61$, is
\begin{center}
{\small\ttfamily 0123456789ABCDEFGHIJKLMNOPQRSTUVWXYZabcdefghijklmnopqrstuvwxyz}
\end{center}
Concatenate the following 436 characters as digits $e_0,\ldots,e_{435}$;
put $\mathcal A=\sum_{j=0}^{435}e_j62^{435-j}$.\par
\begingroup\small\ttfamily
BxEnhvy6NipzkWmROo0GsA3uK3twWLxscLq20UTumy5Cdu461LwD6g7RDNzF8Y\par
iYDFTpB8ehAypd7d6OsOL4APamnJeNXQypUezCmW9n4oXP6XfgmiWuAWZ8ZJ7U\par
OP7gu4kVdlGZqZHULQMar1dXNeJW4i7oNVI8usp8BlovjjZbwPCovZI328fRUq\par
4q92R81s5GjWQaQfGhkq9UYMK3rRMaBmsOIBkJqc2cj2VRIZ6mjzkBIOkW9SjZ\par
FzeA6J5u4j3sBL3SdgwXzM9mHvQivu8EJ1TBTpRhHDqZIbmYgwQ8VqTAFotl6F\par
We1YC554GKrVHsA8oUGKSA6F6k5KREBSD0Tgd7bFALY84LLbpw574BQyvdoHrl\par
0Oe35uwDf15DpT6VJWZBjSbHPpPAINYTte0eUxnTFTKBlGj8Am6sb1gLldh8WZ\par
8P\par
\endgroup

\endgroup


\section{Reproduction and artifact map}

From the extracted package root, run
\begin{lstlisting}
python3 -B verify_all.py
python3 -O -B verify_all.py
\end{lstlisting}
The programs require Python 3.10 or later and its standard library. They
do not search for roots. An evaluator accepts an explicitly provided
integer tuple in the order declared by its equation file. Exact evaluation
may produce large integers; a modular diagnostic is not a proof of an
integer zero. The paper can be rebuilt with a normal TeX Live installation:
\begin{lstlisting}
cd report
pdflatex -interaction=nonstopmode -halt-on-error main.tex
pdflatex -interaction=nonstopmode -halt-on-error main.tex
\end{lstlisting}
Mathematical data are deterministic. PDF bytes can vary with TeX version,
font metadata, or timestamps. The release manifest records the actual
delivered files and should be regenerated after an intentional edit.

\begin{center}
\small
\begin{tabular}{@{}p{0.30\linewidth}p{0.64\linewidth}@{}}
\toprule
Path&Contents\\\midrule
\code{report/}&This paper, its editable LaTeX, printed residuals, and metrics.\\
\code{direct208/}&New literal quartic, expansion, constructor producer, projection
and optimization certificates, exact source audit, and evaluator.\\
\code{positive190/}&Ten-variable graph, whole coefficient payloads and reconstruction,
complete new proof, selected parent proofs, and independent finite checker.\\
\code{affine209/}&Three isolated affine-horizon profiles, parent coefficients,
source commutation certificates, and independent checks.\\
\code{reviews/}&Separately assigned internal mathematical and algebraic reviews;
these are AI-assisted reviews, not external human peer review.\\
\code{dependencies/}&Unchanged incoming Pell Growth 210 archive.\\
\code{verify_all.py}&Ordered, offline verification of the delivered components.\\
\bottomrule
\end{tabular}
\end{center}
\end{document}
